% Tabla: Capas de protección contra lecturas erróneas — v0.17.0
% Arquitectura defensiva del nodo AMI
% Fecha: 2026-03-04

\begin{table}[htbp]
\centering
\caption{Arquitectura de protección multicapa contra lecturas erróneas — v0.17.0}
\label{tab:protection-layers-v017}
\begin{tabular}{c l p{7cm}}
\toprule
\textbf{Capa} & \textbf{Componente} & \textbf{Descripción} \\
\midrule
L1 & \texttt{meter\_read\_all()}
   & Cada lectura OBIS que falla deja su campo en 0 y su bit
     en \texttt{field\_mask} sin activar. Solo lecturas exitosas
     encienden el bit correspondiente. \\
\addlinespace
L2 & \texttt{MIN\_READ\_PERCENT}
   & Si menos del 50\,\% de los códigos OBIS objetivo fueron
     leídos, la lectura completa se marca como inválida
     (\texttt{valid = false}). \\
\addlinespace
L3 & \texttt{readings\_sanity\_check()}
   & Verifica rangos físicos: tensión $\in [50, 500]$\,V,
     frecuencia $\in [40, 70]$\,Hz, y presencia de al menos
     tensión o frecuencia en \texttt{field\_mask}. \\
\addlinespace
L4 & \texttt{THRESH\_CHECK} con \texttt{bit\_idx}
   & Antes de enviar cada campo al servidor LwM2M, verifica
     que su bit esté activo en \texttt{field\_mask}. Campos
     no leídos se omiten (\texttt{skipped++}). \\
\addlinespace
L5 & \texttt{consecutive\_meter\_failures}
   & Contador en \texttt{main.c}: tras 5 fallos consecutivos
     del medidor, se emite un log crítico. Se reinicia
     al primer éxito con mensaje de recuperación. \\
\addlinespace
L6 & \texttt{last\_good} por campo
   & Solo los campos con bit activo en \texttt{field\_mask}
     actualizan su valor en la caché \texttt{last\_good}.
     Campos fallidos conservan su último valor real. \\
\bottomrule
\end{tabular}
\end{table}
